\chapter{АНАЛИЗ ЗАДАЧИ И ПОДХОДОВ К ЕЁ РЕШЕНИЮ}

\section{Постановка задачи управления в игре <<Змейка>>}

В этой задаче нужно построить агента, который управляет змейкой на дискретном поле и старается набрать как можно больше очков. На каждом шаге агент выбирает действие по текущему состоянию игры, после чего среда возвращает новое состояние и награду. Поэтому задачу удобно рассматривать как последовательное принятие решений.

Особенность игры в том, что хороший ход не всегда даёт пользу сразу. Действие может выглядеть безопасным сейчас, но через несколько шагов привести либо к пище, либо к столкновению. Из-за этого агенту нужно учитывать не только ближайшую ситуацию, но и накопленный опыт.

\section{Особенности обучения с подкреплением для дискретной среды}

Обучение с подкреплением подходит для задач, где агент сам ищет удачную стратегию, взаимодействуя со средой и получая награды за свои действия. В отличие от обучения с учителем, здесь нет набора правильных ответов. Агент учится на том, к чему приводят его решения.

На шаге $t$ агент видит состояние $s_t$, выбирает действие $a_t$, получает награду $r_t$ и переходит в состояние $s_{t+1}$. Для игры <<Змейка>> такой подход удобен, потому что:
\begin{itemize}
    \item игра имеет дискретную структуру;
    \item число действий невелико;
    \item результат легко оценивать по счёту;
    \item поведение агента можно наглядно показать на экране.
\end{itemize}

Поэтому <<Змейка>> хорошо подходит для демонстрации базовых идей reinforcement learning на реальном проекте.

\section{Сравнение возможных подходов к решению}

Управлять змейкой можно по-разному. Самый простой вариант --- набор правил и эвристик. В этом случае разработчик заранее прописывает, как агент должен реагировать на опасность и как двигаться к пище. Такой подход легко начать, но он плохо переносится на более сложные игровые ситуации.

Другой вариант --- использовать алгоритмы поиска пути. Они могут строить маршрут к цели, но требуют отдельного планирования и хуже показывают сам процесс обучения на опыте.

Третий вариант --- методы обучения с подкреплением, например табличный Q-learning или более сложные нейросетевые методы, такие как DQN. В нашем проекте состояние уже сведено к компактному набору признаков, поэтому можно применять табличный метод без лишнего усложнения.

\begin{table}[H]
    \centering
    \caption{Сравнение подходов к управлению змейкой}
    \begin{tabularx}{\textwidth}{>{\raggedright\arraybackslash}p{0.22\textwidth} >{\raggedright\arraybackslash}p{0.24\textwidth} X}
        \toprule
        Подход & Преимущества & Ограничения \\
        \midrule
        Правила и эвристики & Простая реализация, предсказуемое поведение & Сложно покрыть все игровые ситуации, слабая адаптивность \\
        Поиск траектории & Явный контроль маршрута, хорош для детерминированных сценариев & Требует дополнительного планирования на каждом шаге, хуже отражает накопление опыта \\
        Табличный Q-learning & Интерпретируемость, простота обновления, естественная связь с дискретной средой & Ограниченная масштабируемость при росте пространства состояний \\
        DQN и другие глубокие методы & Хорошо работают на более сложных наблюдениях и больших пространствах состояний & Требуют больше данных, вычислений и настройки \\
        \bottomrule
    \end{tabularx}
\end{table}

\section{Обоснование выбора табличного Q-learning}

В этой работе выбран табличный Q-learning. На это есть две основные причины. Во-первых, он хорошо подходит для учебного проекта, потому что его легко объяснять и анализировать. Во-вторых, в проекте используется вектор состояния из 11 бинарных признаков, а значит Q-значения ещё можно хранить в таблице.

Такой выбор даёт хороший баланс между простотой и пользой. С одной стороны, метод не перегружает проект лишней сложностью. С другой стороны, он позволяет получить реальный обучаемый агент и показать, как меняется его поведение по мере накопления опыта.
